\chapter{The Operational Semantics of THOR}

Although the pure \lambdacalculus{} is powerful enough to compute all
computable functions, it is not always a convenient medium in which to program.
If one wishes to deal with data objects (integers, for example) a suitable
encoding in terms of \(\lambda\)-expressions must be found. Rather than use such
encodings for common data types, it is more convenient (in both human and
machine terms) to extend the \lambdacalculus{} to include common data types and
rewrite rules for manipulating these types. Such an extension is commonly called
an ``enriched'' or ``applied'' \lambdacalculus.\footnote{RED2 directly supports
the character, integer, floating point, and symbol data types. Special
\(\lambda\)-expression encodings are used to represent structures, which are
compound data types.}

This chapter presents such an extension, The Head Order Reduction Language
(THOR). THOR is a conservative extension of the pure \lambdacalculus, which
means that all expressions reducible in the pure \lambdacalculus{} are
reducible in THOR and all expressions that are not reducible in the pure
\lambdacalculus{} are not reducible in THOR. Like the pure \lambdacalculus,
THOR is untyped.

THOR is the assembly language for RED2. Compared to other assembly languages,
THOR is very high level in nature. THOR is a small but fairly complete
functional programming language, albeit a rather old fashioned one when compared
to ``modern'' functional languages such as Haskell~[26]. THOR is designed so
that it can be used either ``as is'' or as a compilation target for other
functional languages.

This chapter presents a formal semantics for THOR, using Plotkin's structural
approach~[44]. The meaning of each construct in THOR is given by rules which
describe the construct's behavior in terms of abstract machine configurations.
These rules are sometimes rather tedious, but their use is justified because of
the semantic complexities introduced by partial evaluation and incremental
reduction.

\section{The Syntax of Thor}

The domain of THOR is the set of well-formed expressions defined by the
following formulae:
\[
\begin{array}{rcl}
\langle expression\rangle &::=& \langle exp\rangle^+\\
\langle exp\rangle &::=& (\langle expression\rangle) \mid
  \langle abstraction\rangle \mid \langle constant\rangle \mid
  \langle structure\rangle\\
\langle abstraction\rangle &::=&
  \THOR{(LAMBDA (}\langle symbol\rangle^+\THOR{) }\langle expression\rangle\THOR{)}
  \mid \THOR{(LETREC (}\langle binding\text{-}pair\rangle^+\THOR{) }
  \langle expression\rangle\THOR{)}\\
\langle binding\text{-}pair\rangle &::=&
  \THOR{(}\langle symbol\rangle\ \langle exp\rangle\THOR{)}\\
\langle constant\rangle &::=&
  \langle character\rangle \mid \langle integer\rangle \mid
  \langle float\rangle \mid \langle symbol\rangle\\
\langle structure\rangle &::=&
  \THOR{\{}\langle symbol\rangle\ \langle exp\rangle^*\THOR{\}}.
\end{array}
\]
A superscript \(+\) means one or more occurrences and a superscript \(*\) means
zero or more occurrences. A lexical description of each constant type is
omitted; when a constant is presented, its type will be obvious.

The symbol which is the first component of a structure is called its tag. The
tag can loosely be thought of as a type identifier; however, this is a bit
ambiguous because THOR is an untyped language.

There are also metalinguistic mechanisms for associating symbols with
expressions and defining structure types. The association mechanism is analogous
to Scheme's~[15] \THOR{define} special form and is written
\[
s = e
\]
where \(s\) is a symbol and \(e\) is a well-formed expression. All of these
associations together form a definition context within which expressions are
reduced. The mechanism for defining structure types is similar to CommonLisp's
[31] \THOR{defstruct} special form, and is written
\[
s \mapsto s_1\ldots s_n,
\]
where \(s\) is the structure's tag and the \(s_i\) are symbols used to refer to
components of the structure. Such a structure definition generates a set of
\(n\) structure accessors. These mechanisms are only available at the RED2 top
level and are not a part of THOR per se.

The syntax of THOR differs in one major respect from that of the pure
\lambdacalculus{} presented in Chapter 2 - all expressions in THOR are
closed.\footnote{This notion of closed is somewhat different from the
\lambdacalculus's. In the \lambdacalculus{} the meaning of a closed expression
is completely independent of the context in which it occurs. This property does
not necessarily hold for THOR because of the value of symbolic constants depends
on the current definition context.} THOR replaces the identifiers of the pure
\lambdacalculus{} with symbols. Symbols that are within the scope of a
\THOR{LAMBDA}-binding of the same name are considered to be variables; all other
symbols are considered to be symbolic constants. It may appear that this results
in some loss of expressive power, but that is really not the case. Through the
use of \(=\), symbolic constants perform the same role as free variables (i.e.,
they are identifiers for some expression known only in a larger scope).

\section{Notation}

\subsection{Configurations}

The semantics of THOR are presented in terms of configurations, which are
5-tuples describing the state of an abstract machine which reduces THOR
expressions. Configuration tuples are of the form
\[
\cfg{\delta}{\phi}{q}{\rho}{e}
\]
where \(\delta\) is a mapping from symbols to expressions; \(\phi\) is the count
of new \(\lambda\)'s introduced by \(\eta\)-extension; \(q\) is the number of
allowed contractions remaining; \(\rho\) is the redex store; and \(e\) is the
THOR construct being reduced.

The \(\delta\) mapping is the definition context created through use of \(=\) to
associate symbols to expressions; its value is supplied as part of every
reduction sequence's initial configuration and remains invariant during
reduction. The initial value of \(\phi\) is always zero, and \(\rho\) is always
empty (\(\Omega\)) in the initial configuration. The number of allowed
contractions, \(q\), is supplied by the user as part of the initial
configuration and may decrease in value during reduction. The initial expression
\(e\) is supplied by the user as well.

The mappings \(\delta\) and \(\rho\) are partial functions which are made
complete by extending them to return the value bottom (\(\bot\)) everywhere they
are undefined. If \(\delta\) is applied to a symbol which does not have an
associated definition, \(\bot\) is returned. Likewise, if the index supplied to
\(\rho\) exceeds the number of entries in the redex store, \(\bot\) is returned.

The number of contractions initially allowed by the user, \(q\), shall be
referred to as the quantum. When this number is exhausted (i.e., \(q=0\)), I
shall say the quantum has \emph{expired}. (Quantum is an operating systems term
for the amount of time a process is allowed to run before it is swapped out of
the CPU by the scheduler. It seems an appropriate name for the remaining
contraction count.)

\subsection{Reduction Rules}

Reduction rules describe allowable transitions from one configuration to
another. For instance, the rule
\[
\frac{\cfg{\delta}{\phi}{q}{\rho}{x}\redsto
      \cfg{\delta}{\phi}{q}{\rho}{x}}
     {x\in \Int}
\]
specifies that integers reduce to themselves. The \(\redsto\) symbol is
pronounced as ``reduces to.'' A reduction rule consists of a conclusion (the
reduction above the line), zero or more premises (any reductions below the
line), and zero or more application conditions (also below the line). The
conclusion is valid if all premises are valid and all application conditions are
satisfied.

The reduction rules have both a proof theoretic reading and an operational
reading. Under the proof theoretic reading, the rules defining a construct are
used as proof rules to justify instances of the \(\redsto\) relation. A proof
for a reduction is a tree whose structure closely resembles the structure of the
expression being reduced. Under this view, the result of reducing an expression
is computed by constructing such a proof tree.

\begin{figure}
\[
\begin{array}{rcl}
T[\THOR{(LAMBDA }(x_1\ldots x_m)\ e\THOR{)},s]
  &=& \lambda^m T[e,x_1:\cdots:x_m:s]\\
T[\THOR{(LETREC }((x_1\ e_1)\ldots(x_m\ e_m))\ e\THOR{)},s]
  &=& \THOR{(LETREC }(T[e_1,s']\ldots T[e_m,s'])\ T[e,s']\THOR{)}\\
&& \text{where } s'=x_1:\cdots:x_m:s\\
T[\THOR{\{}x\ e_1\ldots e_n\THOR{\}},s]
  &=& \THOR{\{}x\ T[e_1,s]\ldots T[e_n,s]\THOR{\}}\\
T[(e_1e_2),s] &=& (T[e_1,s]\ T[e_2,s])\\
T[x,s] &=& index(x,s,0)
\end{array}
\]
where
\[
\begin{array}{rcl}
index(x,x:s,i)&=&\THOR{(VAR }i\THOR{)}\\
index(x,y:s,i)&=&index(x,s,i+1),\quad x\ne y\\
index(x,\Omega,i)&=&x.
\end{array}
\]
\caption{Translation from THOR to the constructs used in the semantic rules.
The first argument to \(T\) is a THOR expression, and the second argument is the
expression's scope list. The colon (:) is an infix sequence constructor and
\(\Omega\) is the empty sequence.}
\end{figure}

Under the operational reading, the reduction rules specify an abstract
interpreter that initially has a configuration \(\cfg{\delta}{0}{q}{\Omega}{e}\)
and produces a final configuration \(\cfg{\delta}{\phi}{q'}{\rho}{e'}\) such
that
\[
\cfg{\delta}{0}{q}{\Omega}{e}\redsto
\cfg{\delta}{\phi}{q'}{\rho}{e'}
\]
is a valid computation. If the reduction rules correctly implement the
\lambdacalculus, the above implies that under definition context \(\delta\),
\[
e\longrightarrow e'
\]
after \(q-q'\) contractions. The interpreter proceeds by selecting a reduction
rule with a conclusion whose left hand side matches the current configuration
and whose application conditions are satisfied, and then performing the
reductions specified by the premises. This results in a valid configuration on
the right hand side of the conclusion.

The translation from THOR expressions to the constructs used in the reduction
rules is presented in Figure 3.1.

Closures and unbound variables, which are not part of the \lambdacalculus{}
proper but are part of its implementation, are represented by \THOR{(CLOSURE e
p)} and \THOR{(UBV i)} respectively.

\subsection{Miscellaneous}

The set of objects belonging to each data type will be abbreviated as follows:
\(\Char\) is the set of all characters; \(\Int\) is the set of all integers;
\(\Float\) is the set of all floating point numbers; \(\Sym\) is the set of all
symbols; and \(\Struct\) is the set of all structures. Likewise, the classes of
\(\lambda\)-expressions will also be abbreviated: \(\Abstr\) is the set of all
\(\lambda\)-abstractions, \(\App\) is the set of all applications, and \(\Var\)
is the set of variables.

In some of the rules, upper-case letters such as \(E\) will be used to denote an
expression that has already been reduced.

\section{The Pure \(\lambda\)-calculus}

The operational semantics of the pure \lambdacalculus{} is described by six
reduction rules which correspond closely to the final head order reduction rules
presented in Section 2.4.1. The differences between these rule sets can be
accounted for by the abstract machine nature of the operational semantics and
the inclusion of incremental reduction. The meaning of the first rule will be
discussed in some detail in order for the reader to get accustomed to this style
of semantics; later rules will be presented without such detail when it is not
needed.

\subsubsection*{Rule 1 (\(\beta\)-Reduction)}
\[
\frac{
  \cfg{\delta}{\phi}{q_0-1}{\rho'}{e_0}\redsto
  \cfg{\delta}{\phi}{q_1}{\rho'}{e_2}
}{
  \cfg{\delta}{\phi}{q_0}{\rho}{((\lambda e_0)e_1)}\redsto
  \cfg{\delta}{\phi}{q_1}{\rho}{e_2}
}
\quad
\begin{array}{l}
\rho'=\rho+[\THOR{(CLOSURE }e_1\ \rho\THOR{)}]\\
q_0>0
\end{array}
\]
Rule 1 corresponds to \(\beta'''_{ext}\). If there is still quantum remaining
(\(q_0>0\)), the \(\beta\)-redex is contracted. The argument and a pointer to
the current redex store are wrapped in an explicit closure object that is then
added to the redex store; the quantum is decremented. Reduction of \(e_0\)
continues, eventually yielding some expression \(e_2\) after
\(q_0-q_1-1\) contractions have taken place.

If the quantum is exhausted, this rule is not valid; the \(\beta\)-redex will be
treated as an application and interpreted by Rule 2.

Note that in the conclusion, both configurations have the same redex store
\(\rho\), even though \(e_2\) was generated by reducing \(e_0\) in the extended
store \(\rho'\). This is correct, because after \(e_0\) is reduced the extension
is no longer necessary. This implies that an implementation of THOR is free to
discard the extension after \(e_0\) has been reduced, and reclaim the space used
by the extension. Thus, the reduction rules can be used as a basis for proving
the correctness of some aspects of an implementation. Such information will be
exploited in Chapter 4; I will wait until then to point out further
implementation aspects of the reduction rules.

\subsubsection*{Rule 2 (Application)}
\[
\frac{
\cfg{\delta}{\phi}{q_0}{\rho}{e_0}\redsto
\cfg{\delta}{\phi}{q_1}{\rho}{E_0}
\qquad
\cfg{\delta}{\phi}{q_1}{\rho}{(E_0e_1)}\redsto
\cfg{\delta}{\phi}{q_2}{\rho}{e_2}
}{
\cfg{\delta}{\phi}{q_0}{\rho}{(e_0e_1)}\redsto
\cfg{\delta}{\phi}{q_2}{\rho}{e_2}
}
\]
Rule 2 corresponds to \(\Omega'\), but is more restrictive. By specifying the
operator is reduced before the operand, Rule 2 constrains THOR to the normal
order reduction strategy. As was the case with the final head order reduction
rules, it is necessary to impose a partial ordering on these rules so that Rule
1 is applied before Rule 2.

\subsubsection*{Rule 3 (Abstraction)}
\[
\frac{
\cfg{\delta}{\phi+1}{q_0}{\rho'}{e}\redsto
\cfg{\delta}{\phi+1}{q_1}{\rho'}{e'}
}{
\cfg{\delta}{\phi}{q_0}{\rho}{\lambda e}\redsto
\cfg{\delta}{\phi}{q_1}{\rho}{\lambda e'}
}
\quad
\rho'=\rho+[\THOR{(UBV }\phi+1\THOR{)}]
\]
Rule 3 corresponds to \(\eta'''_{ext}\). An \(\eta\)-extension is performed and
the unbound variable index is placed in the redex store.

\subsubsection*{Rule 4 (Variables)}
\[
\cfg{\delta}{\phi}{q}{\rho}{\THOR{(VAR }i\THOR{)}}\redsto
\cfg{\delta}{\phi}{q}{\rho}{\rho(i)}
\]
Rule 4 extracts the value of a variable from the redex store. This value is
either a closure (put in the redex store by Rule 1) or an unbound variable index
(put in by Rule 3). This rule corresponds to \(\beta'''_h\).

\subsubsection*{Rule 5 (Unbound Variables)}
\[
\cfg{\delta}{\phi}{q}{\rho}{\THOR{(UBV }i\THOR{)}}\redsto
\cfg{\delta}{\phi}{q}{\rho}{\THOR{(VAR }\phi-i\THOR{)}}.
\]
Rule 5 replaces unbound variable indices with corrected DeBruijn indices, and
corresponds to \(\Phi\).

\subsubsection*{Rule 6 (Closures)}
\[
\frac{
\cfg{\delta}{\phi}{q_0}{\rho_1}{e}\redsto
\cfg{\delta}{\phi}{q_1}{\rho_1}{e'}
}{
\cfg{\delta}{\phi}{q_0}{\rho_0}{\THOR{(CLOSURE }e\ \rho_1\THOR{)}}\redsto
\cfg{\delta}{\phi}{q_1}{\rho_0}{e'}
}
\]
Rule 6 specifies that the value of a closure is the closure's expression \(e\)
reduced in the closure's context \(\rho_1\). This rule does not directly
correspond to any of the final head order reduction rules; it is necessary
because the abstract machine uses explicit closure objects from which \(e\) and
\(\rho_1\) must be extracted.

\section{Adding in Data Objects}

There are two classes of data objects in THOR: the constants, which are atomic
objects; and structures, which are composite objects built from other structures
or constants.

\subsection{Constants}

With the exception of symbols, the constants reduce to themselves:

\subsubsection*{Rule 7 (Constants)}
\[
\frac{\cfg{\delta}{\phi}{q}{\rho}{x}\redsto
      \cfg{\delta}{\phi}{q}{\rho}{x}}
     {x\in \Int\cup\Float\cup\Char}
\]
If a symbol does not have a definition or the quantum has expired, the symbol is
treated as a symbolic constant which reduces to itself:

\subsubsection*{Rule 8 (Symbolic Constants)}
\[
\frac{\cfg{\delta}{\phi}{q}{\rho}{s}\redsto
      \cfg{\delta}{\phi}{q}{\rho}{s}}
     {s\in\Sym\qquad \delta(s)=\bot\vee q=0}
\]
If a symbol has an expression associated with it, the symbol reduces to that
expression if there is quantum left:

\subsubsection*{Rule 9 (Definitions)}
\[
\frac{
\cfg{\delta}{\phi}{q_0-1}{\rho}{\delta(s)}\redsto
\cfg{\delta}{\phi}{q_1}{\rho}{e}
}{
\cfg{\delta}{\phi}{q_0}{\rho}{s}\redsto
\cfg{\delta}{\phi}{q_1}{\rho}{e}
}
\quad
s\in\Sym,\ q_0>0
\]
Note that it costs one contraction to replace a symbol by its definition.

\subsection{Primitive Operations on Constants}

THOR and RED2 provide a variety of primitive built-in operations for the
constants. They fall into three groups: unary primitives, binary primitives, and
type predicates. All three groups of operators are strict; that is, the correct
number of arguments must be reduced before an operator is applied. Except for
the type of operands the primitives take, the reduction rules for each set of
constants are the same. Therefore, I will only present rules for the integers:

\subsubsection*{Rule 10 (Unary Operator)}
\[
\frac{\cfg{\delta}{\phi}{q}{\rho}{(@\ n)}\redsto
      \cfg{\delta}{\phi}{q-1}{\rho}{e}}
     {q>0\qquad n\in\Int\qquad e=@n}
\]

\subsubsection*{Rule 11 (Binary Operator)}
\[
\frac{\cfg{\delta}{\phi}{q}{\rho}{(@\ n_0\ n_1)}\redsto
      \cfg{\delta}{\phi}{q-1}{\rho}{e}}
     {q>0\qquad n_0,n_1\in\Int\qquad e=n_0@n_1}
\]

\subsubsection*{Rule 12 (Type Predicate - True Case)}
\[
\frac{\cfg{\delta}{\phi}{q}{\rho}{\THOR{(INTEGER? }n\THOR{)}}\redsto
      \cfg{\delta}{\phi}{q-1}{\rho}{\THOR{TRUE}}}
     {q>0\qquad n\in\Int}
\]

\subsubsection*{Rule 13 (Type Predicate - False Case)}
\[
\frac{\cfg{\delta}{\phi}{q}{\rho}{\THOR{(INTEGER? }E\THOR{)}}\redsto
      \cfg{\delta}{\phi}{q-1}{\rho}{\THOR{FALSE}}}
     {q>0\qquad E\notin \Int\cup\App\cup\Var}
\]
The unary and binary operators behave as one would expect, but the type
predicates hold a surprise. Rule 12 specifies that \THOR{(INTEGER? n)} is
\THOR{TRUE} when \(n\) is an integer, but Rule 13 does not specify the result is
\THOR{FALSE} in all other cases, as one might expect. Rule 13 specifies the
result is \THOR{FALSE} only if \(E\) is either a non-integer constant, a
structure, or an abstraction. Therefore, if \(E\) is an unbound variable or an
irreducible application, neither rule is valid.

The semantics for type predicates is specified this way to accommodate partial
evaluation. THOR's policy towards partial evaluation is to perform contractions
only whenever there is enough information available to do so. In the case that
\(E\) is a variable which was not bound by \(\beta\)-reduction, there is no
information about the variable, so a type predicate cannot be contracted. If
\(E\) is an irreducible application, it is also the case that no information is
available. This situation can occur when a symbol in the irreducible expression
has not been associated with a definition; for example,
\[
\THOR{(INTEGER? (FOO x))}
\]
where \THOR{FOO} has not yet been defined.

If \(E\) is an abstraction, there is enough information for type predicates to
make a decision. It is not possible for an argument that is an abstraction to
become anything other than an abstraction; so for example
\[
\THOR{(INTEGER? (LAMBDA (X) e))}
\]
can safely contract to \THOR{FALSE} no matter what \(e\) is.

\subsection{Structures}

Structures are tagged, lazy~[21] composite data objects. Lazy objects in THOR
shield their components from reduction - the components are not reduced when
the structure is created nor when the structure itself is reduced. Components
can only be reduced after they have been extracted from a structure. This is not
to say that the components of a structure are unaffected by reduction; on the
contrary, it is important that their free variables participate in
\(\beta\)-substitution. For example, the expression
\[
\THOR{(LAMBDA (X Y) \{PAIR (FOO X) (* Y 3)\}) 5}
\]
reduces to
\[
\THOR{(LAMBDA (Y) \{PAIR (FOO 5) (* Y 3)\})}.
\]
The components of the \THOR{PAIR} structure will not be reduced until they have
been extracted by an accessor.

The semantics of structures can easily be captured by manipulating \(q\) to
prevent reduction of the component expressions:

\subsubsection*{Rule 14 (Structures)}
\[
\frac{
\cfg{\delta}{\phi}{0}{\rho}{e_i}\redsto
\cfg{\delta}{\phi}{0}{\rho}{e'_i},\ \forall i\in\{0\ldots n\}
}{
\cfg{\delta}{\phi}{q}{\rho}{\{s\ e_0\ldots e_n\}}\redsto
\cfg{\delta}{\phi}{q}{\rho}{\{s\ e'_0\ldots e'_n\}}
}
\]
When a structure is reduced, the number of allowed contractions is temporarily
set to zero. This allows \(\beta\)-substitution of free variables to take place
but prevents redexes from being contracted.

The components of a structure can be extracted by structure accessors, and a
structure's tag can be extracted by the built-in \THOR{TAG} operator:

\subsubsection*{Rule 15 (Structure Accessors)}
\[
\frac{
\cfg{\delta}{\phi}{q_0-1}{\rho}{e_i}\redsto
\cfg{\delta}{\phi}{q_1}{\rho}{e'}
}{
\cfg{\delta}{\phi}{q_0}{\rho}{(s_i\ \{s\ e_0\ldots e_n\})}\redsto
\cfg{\delta}{\phi}{q_1}{\rho}{e'}
}
\quad q_0>0,\ 0\le i\le n
\]

\subsubsection*{Rule 16 (TAG)}
\[
\frac{\cfg{\delta}{\phi}{q}{\rho}{\THOR{(TAG }\{s\ e_0\ldots e_i\}\THOR{)}}\redsto
      \cfg{\delta}{\phi}{q-1}{\rho}{s}}
     {q>0}
\]
There is also a type predicate for structures, which behaves similarly to the
type predicates for constants.

THOR has one predefined type of structure, the \THOR{PAIR}, from which lists are
constructed. The \THOR{PAIR} structure is defined by
\[
\THOR{PAIR := CAR CDR}.
\]
For convenience, a special syntax will be used for lists:
\[
\THOR{[}e_1\mid e_2\THOR{]}
\]
translates into
\[
\THOR{\{PAIR }e_1\ e_2\THOR{\}}.
\]
The symbol \THOR{NIL} is taken to be the empty list, which may also be written
as \THOR{[]}. Multiple item lists which end with the symbol \THOR{NIL} can be
written as
\[
\THOR{[}e_1\ldots e_n\THOR{]},
\]
which translates into nested \THOR{PAIR} structures. The symbols \THOR{CAR} and
\THOR{CDR} can be used to select the first and second components of pairs,
respectively.

\section{Logical Primitives}

In THOR, boolean values are represented by the symbolic constants \THOR{TRUE}
and \THOR{FALSE}. The logical primitives are \THOR{NOT}, \THOR{AND} and
\THOR{OR}.

The \THOR{NOT} is a strict unary operator with simple semantics:

\subsubsection*{Rule 17 (NOT - Case 1)}
\[
\frac{\cfg{\delta}{\phi}{q}{\rho}{\THOR{(NOT TRUE)}}\redsto
      \cfg{\delta}{\phi}{q-1}{\rho}{\THOR{FALSE}}}
     {q>0}
\]

\subsubsection*{Rule 18 (NOT - Case 2)}
\[
\frac{\cfg{\delta}{\phi}{q}{\rho}{\THOR{(NOT FALSE)}}\redsto
      \cfg{\delta}{\phi}{q-1}{\rho}{\THOR{TRUE}}}
     {q>0}
\]
The \THOR{AND} and \THOR{OR} can take any number of arguments, and their
semantics are more complicated. Since these two primitives are symmetrical, only
the semantics for \THOR{AND} will be presented. The simple case occurs when
there is only one argument, which reduces to \THOR{TRUE}:

\subsubsection*{Rule 19 (AND - Case 1)}
\[
\frac{\cfg{\delta}{\phi}{q}{\rho}{\THOR{(AND TRUE)}}\redsto
      \cfg{\delta}{\phi}{q-1}{\rho}{\THOR{TRUE}}}
     {q>0}
\]
When there is more than one argument, \THOR{AND} reduces them one at a time from
left to right. If an argument reduces to \THOR{TRUE}, it is discarded, thereby
reducing the number of arguments by one:

\subsubsection*{Rule 20 (AND - Case 2)}
\[
\frac{
\cfg{\delta}{\phi}{q_0-1}{\rho}{\THOR{(AND }E_0\ldots E_i\ e_i\ldots e_n\THOR{)}}\redsto
\cfg{\delta}{\phi}{q_1}{\rho}{e}
}{
\cfg{\delta}{\phi}{q_0}{\rho}{\THOR{(AND }E_0\ldots E_i\ \THOR{TRUE}\ e_i\ldots e_n\THOR{)}}\redsto
\cfg{\delta}{\phi}{q_1}{\rho}{e}
}
\]
where \(E_j\notin\{\THOR{TRUE},\THOR{FALSE}\}\), \(q_0,n>0\), and \(i\ge0\).
If all the arguments reduce to \THOR{TRUE}, then eventually Rule 19 can be used.
No information is provided by arguments that do not reduce to either
\THOR{TRUE} or \THOR{FALSE}, so these arguments are left alone.

Finally, if any of the arguments reduce to \THOR{FALSE}, the entire \THOR{AND}
construct reduces to \THOR{FALSE}:

\subsubsection*{Rule 21 (AND - Case 3)}
\[
\frac{\cfg{\delta}{\phi}{q}{\rho}{\THOR{(AND }E_0\ldots E_i\ \THOR{FALSE}\ e_i\ldots e_n\THOR{)}}\redsto
      \cfg{\delta}{\phi}{q-1}{\rho}{\THOR{FALSE}}}
     {E_j\notin\{\THOR{TRUE},\THOR{FALSE}\},\ q>0,\ n>0}
\]
As was the case with type predicates, partial evaluation has a strong influence
on \THOR{AND}'s semantics. A decision to reduce the \THOR{AND} construct to
either \THOR{TRUE} or \THOR{FALSE} is made only if enough information is
available. Here are a few examples that illustrate the semantics of the
\THOR{AND}:
\begin{verbatim}
(AND (< 3 4) (ODD? 5)) -> (AND TRUE (ODD? 5))
                         -> (AND (ODD? 5))
                         -> (AND TRUE)
                         -> TRUE

(AND (> 3 4) (ODD? 5)) -> (AND FALSE (ODD? 5))
                         -> FALSE

(LAMBDA (X) (AND (> X 4) (ODD? 5)))
  -> (LAMBDA (X) (AND (> X 4) TRUE))
  -> (LAMBDA (X) (AND (> X 4)))

(LAMBDA (X) (AND (> X 4) (EVEN? 5)))
  -> (LAMBDA (X) (AND (> X 4) FALSE))
  -> (LAMBDA (X) FALSE)
\end{verbatim}

\section{The Conditional Expression}

There is only one conditional primitive in THOR: the \THOR{IF}. The \THOR{IF}
is a ternary operator that is strict only in its first operand. The first
operand is called the condition, the second is the true branch, and the third is
the false branch.

If the condition is \THOR{TRUE}, the \THOR{IF} construct reduces to the true
branch:

\subsubsection*{Rule 22 (IF - True Case)}
\[
\frac{
\cfg{\delta}{\phi}{q_0-1}{\rho}{e_t}\redsto
\cfg{\delta}{\phi}{q_1}{\rho}{e'}
}{
\cfg{\delta}{\phi}{q_0}{\rho}{\THOR{(IF TRUE }e_t\ e_f\THOR{)}}\redsto
\cfg{\delta}{\phi}{q_1}{\rho}{e'}
}
\quad q_0>0
\]

If the condition is \THOR{FALSE}, the \THOR{IF} construct reduces to the false
branch:

\subsubsection*{Rule 23 (IF - False Case)}
\[
\frac{
\cfg{\delta}{\phi}{q_0-1}{\rho}{e_f}\redsto
\cfg{\delta}{\phi}{q_1}{\rho}{e'}
}{
\cfg{\delta}{\phi}{q_0}{\rho}{\THOR{(IF FALSE }e_t\ e_f\THOR{)}}\redsto
\cfg{\delta}{\phi}{q_1}{\rho}{e'}
}
\quad q_0>0
\]

If the condition is not \THOR{TRUE} or \THOR{FALSE}, then neither of the two
branches are reduced:

\subsubsection*{Rule 24 (IF - Non-boolean Case)}
\[
\frac{
\cfg{\delta}{\phi}{q_0}{\rho}{e_c}\redsto
\cfg{\delta}{\phi}{q_1}{\rho}{e'_c}
}{
\cfg{\delta}{\phi}{q_0}{\rho}{\THOR{(IF }e_c\ e_t\ e_f\THOR{)}}\redsto
\cfg{\delta}{\phi}{q_1}{\rho}{\THOR{(IF }e'_c\ e_t\ e_f\THOR{)}}
}
\quad e'_c\notin\{\THOR{TRUE},\THOR{FALSE}\}
\]
None of the reduction rules presented prior to Rule 24 truly depend on the
sequential, normal order reduction strategy of RED2. A different reduction
ordering would still produce the same answer, if given sufficient quantum. Rule
24, however, depends on RED2's specified reduction strategy; the condition must
be reduced before either branch is reduced. This violates the first
Church-Rosser Theorem. As long as THOR runs on the RED2 machine, this is rather
inconsequential. If at some future time THOR is implemented in a parallel
reduction environment, this rule might need to be changed.

Rule 24 is designed to prevent unwarranted recursion in the true and false
branches when the condition is not a boolean value. I consider the problem of
controlling recursion more important than satisfying the Church-Rosser property.

\section{Recursion}

There are three forms of recursion in THOR: explicit single recursion using the
\THOR{Y} operator; explicit multiple recursion using the \THOR{LETREC}
abstraction mechanism; and implicit recursion accomplished through symbols that
have associated definitions.

\subsection{The Y Operator}

The \THOR{Y} operator, or fixpoint combinator as it is sometimes called, is used
for creating explicit recursion. It has a simple semantics:

\subsubsection*{Rule 25 (Y Combinator)}
\[
\frac{
\cfg{\delta}{\phi}{q_0-1}{\rho}{(e\ \THOR{(Y }e\THOR{)})}\redsto
\cfg{\delta}{\phi}{q_1}{\rho}{e'}
}{
\cfg{\delta}{\phi}{q_0}{\rho}{\THOR{(Y }e\THOR{)}}\redsto
\cfg{\delta}{\phi}{q_1}{\rho}{e'}
}
\quad q_0>0
\]
The \THOR{Y} can be used to create recursive functions and represent infinite
data structures. An expression for computing the factorial of an integer is
presented as an example of recursion using the \THOR{Y} operator:
\begin{verbatim}
(Y (lambda (fact n) (if (= n 0) 1 (* n (fact (1- n))))))
\end{verbatim}
Note that \THOR{fact} is a variable, which gets bound to the entire recursive
expression during reduction.

The fixpoint combinator could also be implemented by any of several
\(\lambda\)-expressions, including
\[
(\lambda x.((\lambda y.x(yy))(\lambda y.x(yy))))
\]
which is due to Curry~[25]. The advantage of using a \(\lambda\)-expression for
\THOR{Y} is that no new reduction mechanisms need be introduced. There are
several disadvantages to using \(\lambda\)-expressions to implement the fixpoint
combinator, however. For one thing, implementing \THOR{Y} directly as a
primitive machine operation is more efficient in both time and space than using
a \(\lambda\)-expression. But more importantly, the \THOR{Y} primitive retains
its identity throughout a reduction sequence.

Retaining identity is important in an incremental reduction system. When a
person steps through a reduction sequence, it is desirable that the output
(reduced) expressions are easily recognized as resulting from the input
expressions. Most of the primitives in THOR, such as the arithmetic primitives,
are local in nature and easily retain their identity. The fixpoint combinator is
global in nature and can affect many parts of an expression, so retaining
identity is more difficult.

A comparison of the reduction sequences of the factorial expression described
earlier illuminate the identity problems associated with recursion. Using a
\(\lambda\)-expression to implement the fixpoint combinator yields a reduction
sequence that gets complicated quickly; if you were shown only the third
expression, would you be able to clearly identify what was going on? Now examine
the reduction sequence for the \THOR{Y} primitive:
\begin{verbatim}
((Y (LAMBDA (FACT N) (IF (= N 0) 1 (* N (FACT (1- N)))))) 3)
  -> ((LAMBDA (N)
        (IF (= N 0)
            1
            (* N ((Y (LAMBDA (FACT N)
                       (IF (= N 0) 1 (* N (FACT (1- N))))))
                  (1- N)))))
      3)
  -> (IF (= 3 0)
         1
         (* 3 ((Y (LAMBDA (FACT N)
                    (IF (= N 0) 1 (* N (FACT (1- N))))))
               (1- 3))))
\end{verbatim}
The latter sequence is easier to comprehend. Use of the \THOR{Y} operator does
not always make recursive reduction sequences easy to follow, but it does make
recursion easier to spot than using a \(\lambda\)-expression.

\subsection{Mutual Recursion with LETREC}

The \THOR{LETREC} abstraction mechanism provides a way to create local mutually
recursive expressions. It is a complicated mechanism, so it will be described in
some detail. As a reminder of its syntax, here again is its form:
\[
\THOR{(LETREC ((}var_1\ exp_1\THOR{) ... (}var_n\ exp_n\THOR{)) }body\THOR{)}
\]
The \THOR{LETREC} allows any \(var_i\) to be correctly referenced inside any of
the \(exp_i\) and the body. As a highly contrived example of mutual recursion,
the following expression determines if an integer is even:
\begin{verbatim}
(lambda (x)
  (letrec ((even? (lambda (n) (if (= n 0) TRUE (odd? (1- n)))))
           (odd?  (lambda (n) (if (= n 0) FALSE (even? (1- n))))))
    (even? x)))
\end{verbatim}

The \THOR{LETREC} can be implemented in the following way. Each of the
\((var_i\ exp_i)\) pairs form a \(\beta\)-redex. When the \THOR{LETREC} is
reduced, the redex store \(\rho\) is updated so that:
\[
\rho'=\rho+[\THOR{(CLOSURE }exp_1\ \rho'\THOR{)},\ldots,
             \THOR{(CLOSURE }exp_n\ \rho'\THOR{)}].
\]
The new \(\rho'\) will be called the recursive context of the \THOR{LETREC}.
Mutually recursive references are allowed in the \(exp_i\) because the redex
store pointer of each closure in \(\rho'\) is such that it includes the entire
recursive context.

The need to accommodate incremental reduction, however, invalidates this simple
approach. The recursive context created by a \THOR{LETREC} may need to exist
throughout the entire reduction of the \THOR{LETREC}'s body, which may take many
contractions. However, the quantum may expire before the body is completely
reduced. Therefore, some sort of mechanism must exist for reconstructing the
recursive context wherever it will be needed again. This can be done by
reconstructing the \THOR{LETREC}'s original form around every one of the
recursive variables that exist in the partially reduced body.

An example of the desired behavior helps clarify the discussion. The expression
\begin{verbatim}
(letrec ((x [1 | y])
         (y [2 | x]))
  x)
\end{verbatim}
describes an infinite list of alternating 1's and 2's. Reducing with a quantum
of 1 allows the \THOR{LETREC} to create the recursive context and have its body,
\THOR{x}, be expanded to \THOR{[1 | y]}. But now the quantum is exhausted, and
all of the recursion variables in the body must be ``wrapped'' in a
reconstructed \THOR{LETREC}. So the expression is now:
\begin{verbatim}
[1 | (letrec ((x [1 | y]) (y [2 | x])) y)]
\end{verbatim}
The cdr of this list is now of a form similar to that of the original
expression, and its reduction sequence unfolds in a similar way. Ignoring for a
moment that the \THOR{PAIR} structure comprising the list prevents the reduction
of its components, another reduction gives
\begin{verbatim}
[1 2] (letrec ((x [1 | y]) (y [2 | x])) x)]
\end{verbatim}
The reduction sequence can continue on in this fashion, ad infinitum.

Two new types of object are introduced to implement the recursive context and
help reconstruct \THOR{LETREC} binding expressions: \THOR{BLOCK}s and
\THOR{REC}s. A \THOR{BLOCK} saves all of the binding expressions so that they
may be recalled, and \THOR{REC}s provides a way to access the binding
expressions stored in a \THOR{BLOCK}.

\subsubsection*{Rule 26 (LETREC - Case 1)}
\[
\frac{
\cfg{\delta}{\phi}{q_0}{\rho'}{e}\redsto
\cfg{\delta}{\phi}{q_1}{\rho'}{e'}
}{
\cfg{\delta}{\phi}{q_0}{\rho}{\THOR{(LETREC }(e_1\ldots e_n)\ e\THOR{)}}\redsto
\cfg{\delta}{\phi}{q_1}{\rho}{e'}
}
\]
\[
\rho'=\rho+[\THOR{(REC }n\ \rho'\ b\THOR{)},\ldots,
            \THOR{(REC }1\ \rho'\ b\THOR{)}],
\qquad b=\THOR{(BLOCK }e_1\ldots e_n\THOR{)},\quad q_0>0.
\]
When a \THOR{LETREC} is reduced, a new \THOR{BLOCK} is created containing all of
the binding expressions and a \THOR{REC} is placed in the redex store for each
binding expression. There is no contraction charge for creating this recursion
context.

If a \THOR{LETREC} is to be reduced after the quantum has expired, the recursive
context is not created. Instead, each of the recursive bindings is treated like
an unbound variable. The binding expressions and the body are allowed to
participate in \(\beta\)-substitution:

\subsubsection*{Rule 27 (LETREC - Case 2)}
\[
\frac{
\cfg{\delta}{\phi+n}{0}{\rho'}{e}\redsto
\cfg{\delta}{\phi+n}{0}{\rho'}{e'}
\qquad
\cfg{\delta}{\phi+n}{0}{\rho'}{e_i}\redsto
\cfg{\delta}{\phi+n}{0}{\rho'}{e'_i}
}{
\cfg{\delta}{\phi}{0}{\rho}{\THOR{(LETREC }(e_1\ldots e_n)\ e\THOR{)}}\redsto
\cfg{\delta}{\phi}{0}{\rho}{\THOR{(LETREC }(e'_1\ldots e'_n)\ e'\THOR{)}}
}
\]
\[
\rho'=\rho+[\THOR{(UBV }\phi+1\THOR{)},\ldots,\THOR{(UBV }\phi+n\THOR{)}].
\]

When a \THOR{REC} is selected from the redex store by a variable, it reduces to
its corresponding binding expression stored in its \THOR{BLOCK}:

\subsubsection*{Rule 28 (Rec - Case 1)}
\[
\frac{
\cfg{\delta}{\phi}{q_0-1}{\rho_1}{e_i}\redsto
\cfg{\delta}{\phi}{q_1}{\rho_1}{e'}
}{
\cfg{\delta}{\phi}{q_0}{\rho_0}{\THOR{(REC }i\ \rho_1\ b\THOR{)}}\redsto
\cfg{\delta}{\phi}{q_1}{\rho_0}{e'}
}
\quad b=\THOR{(BLOCK }e_1\ldots e_n\THOR{)},\ q_0>0
\]

Reducing a \THOR{REC} costs one contraction. If the quantum has expired when a
\THOR{REC} is reduced, the \THOR{LETREC} construct must be reconstructed and
``wrapped'' around the variable the \THOR{REC} represents:

\subsubsection*{Rule 29 (Rec - Case 2)}
\[
\frac{
\cfg{\delta}{\phi+n}{0}{\rho'}{e_j}\redsto
\cfg{\delta}{\phi+n}{0}{\rho'}{e'_j},\ \forall j\in\{1\ldots n\}
}{
\cfg{\delta}{\phi}{0}{\rho_0}{\THOR{(REC }i\ \rho_1\ b\THOR{)}}\redsto
\cfg{\delta}{\phi}{0}{\rho_0}{\THOR{(LETREC }(e'_1\ldots e'_n)\ \THOR{(VAR }i\THOR{))}}
}
\]
\[
\rho'=\rho_0+[\THOR{(UBV }\phi+1\THOR{)},\ldots,\THOR{(UBV }\phi+n\THOR{)}],
\qquad b=\THOR{(BLOCK }e_1\ldots e_n\THOR{)}.
\]

As is the case with the fixpoint operator, mutual recursion can also be
accomplished using special \(\lambda\)-expressions. The reasons for preferring
the \THOR{LETREC} primitive are identical to those given for the \THOR{Y}
primitive. In the case of mutual recursion, the argument is even stronger
because expressions get larger and more complicated than in the single recursion
case.

\subsection{Recursion via Symbol Definitions}

Recursion can be obtained implicitly through the reduction of symbols that have
been associated with expressions by the use of the \(=\) mechanism. This is the
method of recursion most people are familiar with. To contrast this method with
using the \THOR{Y} operator, here is another expression for computing
factorials:
\begin{verbatim}
fact = (lambda (n) (if (= n 0) 1 (* n (fact (1- n)))))
\end{verbatim}
In this expression, \THOR{fact} is not a variable, but a symbolic constant (that
just happens to be on both sides of the definition). When the symbol
\THOR{fact} is encountered as an operator during reduction, it is replaced by
its defining expression. Rule 8 governs this situation.

This method of recursion is generally used in LISP systems~[3]. It is a popular
implementation technique because it allows all global definitions to be mutually
recursive.

\subsection{Controlling Recursion}

If checks are not placed on recursion, it will often consume all of the
available contractions. For example, the nonsense expression \THOR{(Y 1)} will
continue expanding until the quantum is exhausted.\footnote{Although the
resulting expression, \THOR{(1 1 1 ...)}, may not have any meaning to a
particular human observer, it is a legal THOR expression. THOR is filled with
such ``nonsense'' expressions; that makes THOR wonderfully flexible and allows
you to assign whatever meaning you wish to many kinds of expressions.}
Fortunately, most uses of recursion are sensible in nature. Recursion is
generally used for two purposes: implementing self-referential functions and
representing infinite data structures. Each of these uses provides a different
way of controlling recursive expansion.

When recursion is used in self-referential functions, it is usually within the
context of some conditional which prevents recursing infinitely. For example, in
the \THOR{FACT} function defined earlier, the conditional test \THOR{(= N 0)} is
used to stop recursion when \(N\) becomes zero. This is all good, fine, and
obvious - but what about situations involving partial evaluation, where
conditional tests might not reduce to either \THOR{TRUE} or \THOR{FALSE}? I
would like to reduce
\[
\THOR{(lambda (x) (fact x))}
\]
and get back
\[
\THOR{(lambda (x) (if (= x 0) 1 (* x (fact (1- x)))))}
\]
instead of an expression where the \THOR{FACT} in the false branch of the
\THOR{IF} keeps expanding until the quantum is exhausted. It is for this reason
that the semantics of \THOR{IF} specify that neither the true nor false branch
is reduced if the condition does not reduce to \THOR{TRUE} or \THOR{FALSE}. By
not reducing either branch in this example, the recursive expansion of
\THOR{FACT} is prevented.

The charge of one contraction for reducing a \THOR{REC} is also part of the
machinery for controlling recursion. If there were no charge, \THOR{REC}s could
infinitely expand during the \(\beta\)-substitution of the true and false
branches of an \THOR{IF} with a non-boolean condition. Because of the charge,
Rule 29 is used in this situation instead of Rule 28.

The lazy semantics of structures control the recursive expansion of infinite
data structures. For example,
\[
\THOR{(Y (lambda (x) [1 | x]))}
\]
reduces to
\[
\THOR{[1 | (Y (lambda (x) [1 | x]))]},
\]
even if given infinite quantum.

\section{The \(\alpha\)-Equality Predicate}

The final primitive operation to be presented is the \THOR{EQUAL?} predicate,
which tests two expressions for \(\alpha\)-equality. In the \lambdacalculus, two
expressions are \(\alpha\)-equivalent if they are structurally (syntactically)
the same except for their bound variable names.

\THOR{EQUAL?} is a strict binary operator that works by decomposing its
arguments into their atomic constituents and then checking these for syntactic
equality. If all the parts are equal, the result is \THOR{TRUE}; if the parts
are not equal, the result is \THOR{FALSE}; and if some of the parts are unbound
variables, then the result will be a new expression that is a conjunction of
equality tests on the unbound variables. If any of the constituents are
structures, the components of the structures will be reduced before they are
compared. The test for \(\alpha\)-equality may take many contractions to
complete.

The arguments to \THOR{EQUAL?} are first reduced, and then passed on to the
primitive \THOR{EQUAL*}\footnote{\THOR{EQUAL*} is internal to RED2 and is not
directly accessible to users.} which treats the arguments as purely syntactic
objects:

\subsubsection*{Rule 30 (Alpha Equality)}
\[
\frac{
\cfg{\delta}{\phi}{q_0}{\rho}{\THOR{(EQUAL* 0 }E_0\ E_1\THOR{)}}\redsto
\cfg{\delta}{\phi}{q_1}{\rho}{e}
}{
\cfg{\delta}{\phi}{q_0}{\rho}{\THOR{(EQUAL? }E_0\ E_1\THOR{)}}\redsto
\cfg{\delta}{\phi}{q_1}{\rho}{e}
}
\]
The first argument to \THOR{EQUAL*} is a count of the \(\lambda\)'s that have
been encountered so far in the decomposition of \(E_0\) and \(E_1\). This value
will be needed in case the quantum expires before the equality check is
completed.

\THOR{EQUAL*} proceeds to decompose compound arguments into their components,
and check the components for equality:

\subsubsection*{Rule 31 (Equality: Applications)}
\[
\frac{
\cfg{\delta}{\phi}{q_0-1}{\rho}{\THOR{(AND (EQUAL* }n\ E_0\ E_2\THOR{) (EQUAL* }n\ E_1\ E_3\THOR{))}}\redsto
\cfg{\delta}{\phi}{q_1}{\rho}{e}
}{
\cfg{\delta}{\phi}{q_0}{\rho}{\THOR{(EQUAL* }n\ (E_0\ E_1)\ (E_2\ E_3)\THOR{)}}\redsto
\cfg{\delta}{\phi}{q_1}{\rho}{e}
}
\quad q_0>0
\]

\subsubsection*{Rule 32 (Equality: Abstractions)}
\[
\frac{
\cfg{\delta}{\phi}{q_0-1}{\rho}{\THOR{(EQUAL* }n+1\ E_0\ E_1\THOR{)}}\redsto
\cfg{\delta}{\phi}{q_1}{\rho}{e}
}{
\cfg{\delta}{\phi}{q_0}{\rho}{\THOR{(EQUAL* }n\ \lambda.E_0\ \lambda.E_1\THOR{)}}\redsto
\cfg{\delta}{\phi}{q_1}{\rho}{e}
}
\quad q_0>0
\]

\subsubsection*{Rule 33 (Equality: Structures)}
\[
\frac{
\cfg{\delta}{\phi}{q_0-1}{\rho}{\THOR{(AND }(equal_0)\ldots(equal_i)\THOR{)}}\redsto
\cfg{\delta}{\phi}{q_1}{\rho}{e}
}{
\cfg{\delta}{\phi}{q_0}{\rho}{\THOR{(EQUAL* }n\ \{s\ e_0\ldots e_i\}\ \{s\ e'_0\ldots e'_i\}\THOR{)}}\redsto
\cfg{\delta}{\phi}{q_1}{\rho}{e}
}
\quad q_0>0
\]
where \(equal_j=\THOR{(EQUAL? }\lambda^n.e_j\ \lambda^n.e'_j\THOR{)}\).
Rule 32 strips off a \(\lambda\) from the expressions, adding one to the
\(\lambda\) count. Rule 33 decomposes structures, reducing the structure
components before testing them for pairwise equality. Before reducing the
components, however, it is necessary to re-create the proper binding context so
that the DeBruijn indices inside the component expressions are correct. Note
that both structures' tag must be the same and the structures must have the same
number of components.

If the arguments being compared are identical atomic objects, then the result of
\THOR{EQUAL*} is \THOR{TRUE}:

\subsubsection*{Rule 34 (Equality: Identical Atomic Objects)}
\[
\frac{\cfg{\delta}{\phi}{q}{\rho}{\THOR{(EQUAL* }n\ E\ E\THOR{)}}\redsto
      \cfg{\delta}{\phi}{q-1}{\rho}{\THOR{TRUE}}}
     {E\notin\Struct\cup\App\cup\Abstr,\ q>0}
\]
Note that two identical unbound variables are \(\alpha\)-equal; the expression
\[
\THOR{(LAMBDA (X) (EQUAL? X X))}
\]
reduces to \(\THOR{(LAMBDA (X) TRUE)}\).

If one of the arguments is an unbound variable (or if both arguments are
different unbound variables), then a decision about equality cannot be made. If
this is the case, the arguments are wrapped back up in the appropriate number of
\(\lambda\)'s and left alone:

\subsubsection*{Rule 35 (Equality: Variable 1)}
\[
\frac{\cfg{\delta}{\phi}{q}{\rho}{\THOR{(EQUAL* }n\ \THOR{(VAR }i\THOR{)}\ E\THOR{)}}\redsto
      \cfg{\delta}{\phi}{q}{\rho}{\THOR{(EQUAL? }\lambda^n.\THOR{(VAR }i\THOR{)}\ \lambda^n.E\THOR{)}}}
     {i>n}
\]

\subsubsection*{Rule 36 (Equality: Variable 2)}
\[
\frac{\cfg{\delta}{\phi}{q}{\rho}{\THOR{(EQUAL* }n\ E\ \THOR{(VAR }i\THOR{)}\THOR{)}}\redsto
      \cfg{\delta}{\phi}{q}{\rho}{\THOR{(EQUAL? }\lambda^n.E\ \lambda^n.\THOR{(VAR }i\THOR{)}\THOR{)}}}
     {i>n}
\]
Note the stipulation that \(i>n\) in both rules; this means that \THOR{(VAR i)}
must be free with respect to the original expressions being checked for
\(\alpha\)-equality. This condition ensures that
\[
\THOR{(EQUAL? (LAMBDA (X) X) (LAMBDA (Y) 2))}
\]
reduces to \THOR{FALSE}, while
\[
\THOR{(LAMBDA (X) (EQUAL? (LAMBDA (Y) X) (LAMBDA (Y) 2)))}
\]
reduces to itself because there is not enough information to decide whether
\THOR{X} is 2.

If the arguments are not \(\alpha\)-equal, the result is \THOR{FALSE}:

\subsubsection*{Rule 37 (Equality: Unequal Objects)}
\[
\frac{\cfg{\delta}{\phi}{q}{\rho}{\THOR{(EQUAL* }n\ E_0\ E_1\THOR{)}}\redsto
      \cfg{\delta}{\phi}{q-1}{\rho}{\THOR{FALSE}}}
     {E_0\ne E_1,\ q>0}
\]
The condition \(E_0\ne E_1\) does not need to call on some higher authority to
check for inequality\footnote{If such an authority existed, then it could have
been called on to check for equality, saving the need for so many rules.} if a
partial ordering is imposed on Rules 31-37 so that Rule 37 is tried last. If
Rule 37 is ever valid under this ordering, then this condition must be true.

Finally, there is the case of an expired quantum:

\subsubsection*{Rule 38 (Equality: Expired Quantum)}
\[
\frac{\cfg{\delta}{\phi}{0}{\rho}{\THOR{(EQUAL* }n\ E_0\ E_1\THOR{)}}\redsto
      \cfg{\delta}{\phi}{0}{\rho}{\THOR{(EQUAL? }\lambda^n.E_0\ \lambda^n.E_1\THOR{)}}}
     {}
\]
The arguments are wrapped back up in their \(\lambda\)'s, and a new
\THOR{EQUAL?} expression is returned. Here are a few examples to illustrate the
behavior of \THOR{EQUAL?}:
\begin{verbatim}
(equal? 1 1) -> true
(equal? (a b) (a b)) -> true
(lambda (x) (equal? (a 1) (x 1))) -> (lambda (x) (equal? a x))
(equal? (lambda (a) a) (lambda (b) b)) -> true
(equal? (lambda (x) 1) (lambda (x) x)) -> false
(lambda (x) (equal? [1 x 3] [1 2 3]))
  -> (lambda (x) (equal? x 2))
(equal? {node 2 (+ 1 4)} {node 2 5}) -> true
\end{verbatim}
